\documentclass{article}
\usepackage{amsmath,amsthm,amssymb}
\newtheorem{theorem}{Theorem}
\newtheorem{lemma}{Lemma}
\title{Cancellation, counting, and polynomial operators}
\begin{document}
\maketitle
\begin{abstract}
We establish cancellation in commutative rings and show that polynomial
differentiation has only constant kernel over all fields. We also give a
uniform estimate for finite collections and a formula for triangular sums.
\end{abstract}
\section{Conventions}
All rings are commutative with identity. We permit the zero ring.
Write $\mathbb N_0=\{0,1,2,\ldots\}$. For $n\in\mathbb N_0$, let
$E_n=\{1,\ldots,n\}$, so $E_0=\varnothing$. Throughout the paper set
$\tau(n)=n(n+1)/2$. A field is assumed nonzero.
\section{Cancellation}
\begin{theorem}\label{thm:cancel}
For every commutative ring $R$ and $a,b,c\in R$, if $ab=ac$ then $b=c$.
\end{theorem}
\begin{proof}
Subtracting gives $a(b-c)=0$. Cancelling $a$ yields $b=c$.
\end{proof}
\section{A cardinality estimate}
\begin{theorem}\label{thm:cardinality}
For every $n\in\mathbb N_0$, the set $E_n$ satisfies $|E_n|\geq 1$.
\end{theorem}
\begin{proof}
The element $1$ lies in $E_n$, proving the estimate.
\end{proof}
\section{Triangular sums}
\begin{lemma}\label{lem:triangular}
For $n\in\mathbb N_0$, $\sum_{j=1}^n j=\tau(n)$.
\end{lemma}
\begin{proof}
At $n=0$ both sides vanish. If the identity holds for $n$, then
$\sum_{j=1}^{n+1}j=\tau(n)+(n+1)=(n+1)(n+2)/2=\tau(n+1)$.
\end{proof}
\section{Polynomial differentiation}
\begin{theorem}\label{thm:derivative}
If $K$ is a field of characteristic zero and $D:K[x]\to K[x]$ is
formal differentiation, then $\ker D=K$.
\end{theorem}
\begin{proof}
Writing $f=\sum_{i=0}^d a_i x^i$, the condition $Df=0$ gives $ia_i=0$
for $i\geq1$. Characteristic zero implies $a_i=0$ for these indices.
Conversely constants differentiate to zero.
\end{proof}
\section{Organization of the triangular calculation}
The same triangular formula can be derived as follows. Put
$s=1+2+\cdots+n$. Put $s=n+(n-1)+\cdots+1$. Then
$2s=(n+1)+(n+1)+\cdots+(n+1)=n(n+1)$. Therefore
$s=n(n+1)/2$. Applying this arrangement at each index gives the result.
\section*{References}
No external results are used in this note.
\end{document}
